\paragraph{\trPlurisubharmonic{}函数。}\label{sub:plurisubharmonic-functions}
一个\textbf{\trStein{}\trManifold}（Stein manifold）是一个容许\textbf{\trExhausting}（exhausting）——意即\trProper{}且有下界——的\textbf{\trPlurisubharmonic{}函数}（plurisubharmonic function）$\phi\colon X\to\RR$ 的复\trManifold{}$(X,I_X)$。\trPlurisubharmonic{}性是说，若令 $\theta_\phi=-d^c\phi=-d\phi\circ I_X$，则 $\omega_\phi=d\theta_\phi=-dd^c\phi$ 是一个\trKaehler{}\trForm。注意 $dd^c\phi$ 总是 $(1,1)$ 型的；非平凡的条件是它应当在复切线上为负（即 $TX$ 的实 $I_X$-不变二维子空间，带自然\trOrientation）。一个直接的推论是：给定任何\trHolomorphic{}映射 $u\colon S\to X$，函数 $\rho=\phi\circ u$ 满足如 \eqref{eq:rho-laplacian} 中的 $\Delta\rho\ge 0$，从而服从\trMaximumPrinciple。

由于\trPlurisubharmonic{}性条件本质上是凸的，这类函数可以相对容易地加以操作。我们需要这一点的两个初等实例（沿这些路线的更多内容见 [42] 或 [22]）。
\begin{lemma}\label{thm:plurisubharmonic-monotone-recomposition}
设 $\phi$ 是一个\trExhausting{}的\trPlurisubharmonic{}函数，而 $h\colon\RR\to\RR$ 是一个处处 $h'>0$、$h''\ge 0$ 的函数。那么 $h(\phi)$ 仍是一个\trExhausting{}的\trPlurisubharmonic{}函数。
\end{lemma}
\begin{proof}
我们有 $dd^c h(\phi)=h'(\phi)\,dd^c\phi+h''(\phi)\,d\phi\wedge d^c\phi$。在复切线上，第一个加项为负，第二个非正。
\end{proof}
\begin{lemma}\label{thm:plurisubharmonic-modification}
设 $\phi,\psi$ 是 $(X,I_X)$ 上的\trExhausting{}\trPlurisubharmonic{}函数。假设 $\psi$ 比 $\phi$ 增长得更快，意即 $\psi-\phi$ 仍是\trExhausting{}的。取定 $X$ 的某个\trCompact{}子集。那么存在另一个\trExhausting{}的\trPlurisubharmonic{}函数 $\tilde{\psi}$，使得 $\tilde{\psi}$ 在该\trCompact{}子集上等于 $\phi$，而在无穷远处 $\tilde{\psi}=\psi+$（常数）。
\end{lemma}
\begin{proof}
我们选取一个函数 $l$，使得对 $t\ll 0$ 有 $l(t)=0$，对 $t\gg 0$ 有 $l(t)=t$，并且处处 $l''(t)\ge 0$。取 $\xi=\phi+l(\psi-\phi-b)$，其中 $b$ 是某个常数。出于与前面相同的理由，这是\trPlurisubharmonic{}的：$dd^c\xi=(1-l'(\psi-\phi-b))(dd^c\phi)+l'(\psi-\phi-b)(dd^c\psi)+l''(\psi-\phi-b)\,d(\psi-\phi)\wedge d^c(\psi-\phi)$。通过选取足够大的常数 $b$，我们可以保证 $l(\psi-\phi-b)$ 在所选\trCompact{}子集上等于 $0$。
\end{proof}
我们说 $\phi$ 是\textbf{\trFiniteType}（of finite type）的，如果它的\trCriticalPoint{}集是\trCompact{}的。一般而言，判定一个 \trStein{}\trManifold{}是否容许这样一个函数可能是困难的（对四维\trManifold，存在着来自\trDifferential{}结构的微妙\trObstruction{} [63]）。然而，有好几大类情形可以轻易做到。
\begin{example}\label{thm:plurisubharmonic-examples}
（i）设 $\bar{X}$ 是一个\trSmooth{}复\trProjectiveVariety，装备一个\trAmple{}\trLineBundle{}$\scr{L}$。\trAmple{}性意味着，在该\trLineBundle{}上存在一个\trMetric{}$|\cdot|$，其\trCurvature{}是 $\bar{X}$ 上的一个\trKaehler{}\trForm。设 $s$ 是 $\scr{L}$ 的一个\trHolomorphic{}\trSection，使得 $s^{-1}(0)$ 是一个\trNormalCrossings{}\trDivisor。那么
\begin{equation*}
\phi(x)=-\log|s(x)|
\end{equation*}
是 $X=\bar{X}\setminus s^{-1}(0)$ 上一个\trExhausting{}的严格\trPlurisubharmonic{}的\trFiniteType{}函数。相伴的\trKaehler{}\trForm{}$\omega_\phi$ 是先前给出的那个\trKaehler{}\trForm{}在 $X$ 上的限制。\trFiniteType{}性质由一个直截了当的局部计算所证明 [133, Lemma 8]。

（ii）设 $X\subset\CC^N$ 是一个\trSmooth{}仿射\trAlgebraicVariety，$\phi\colon X\to\RR$ 是一个\trExhausting{}的\trPlurisubharmonic{}函数，它是某个实多项式的（限制）。那么 $\phi$ 自动是\trFiniteType{}的，因为它的\trCriticalPoint{}集本身是一个实\trAlgebraicVariety（从而只有有限多个\trConnectedComponent，每个都对应于单一的\trCriticalValue）。

（iii）沿着同一思路，设 $q\colon\CC^N\to\CC$ 是一个在原点处有\trIsolated{}\trCriticalPoint{}的复多项式，其中 $q(0)=0$、$Dq(0)=0$。那么，对一切 $0<|\epsilon|\ll\delta$，\trMilnorFibre
$X=\{x\in\CC^N: q(x)=\epsilon,\ |x|<\delta\}$ 是一个 \trStein{}\trManifold，而 $\phi(x)=(\delta^2-|x|^2)^{-1}$ 是其上一个\trFiniteType{}的\trPlurisubharmonic{}函数。注意，\trMilnorFibre{}一般并非（\trBiholomorphic{}于）\trAlgebraicVariety。
\end{example}
从现在起，取定 $(X,I_X)$，它另外还应配有一个二次复\trVolume{}形式 $\eta_X^2$ 与一个\trFiniteType{}的\trExhausting{}\trPlurisubharmonic{}函数 $\phi$。取 $c$ 大于 $\phi$ 的一切\trCriticalValue，并考虑\trSublevelSet
\begin{equation*}
M=X^{\phi,c}=\phi^{-1}((\infty;c])
\end{equation*}
连同 $\theta_X$、$\omega_X$ 的限制。这是一个带\trContactTypeBoundary{}的\trExactSymplectic\trManifold（而 $I_X$ 的限制在边界附近是\trContact{}型的\trAlmostComplex{}结构）。再配上 $\eta_X^2$ 的限制，我们便有了定义\trFukayaCategory{}$\scr{F}(M)$ 所需的全部数据。
\begin{proposition}\label{thm:fukaya-category-independent}
在相差\trQuasiequivalence{}的意义下，$\scr{F}(M)$ 与 $c$ 和 $\phi$ 的选取无关。
\end{proposition}
\begin{proof}
全部讨论会与\cref{sub:contact-type-boundary-case}中的讨论多少有些类似。第一步是表明，\emph{对固定的 $\phi$，$c$ 的选取无关紧要。}取某个别的 $d>c$。在定义 $\scr{F}(X^{\phi,d})$ 时可以安排得使：对含于较小子集 $M=X^{\phi,c}$ 之内的那些\trExact{}\trLagrange\trSubmanifold，其\trFloerDatum{}与\trPerturbationDatum{}只使用在 $X^{\phi,d}\setminus\mathrm{int}(X^{\phi,c})$ 的一个邻域上等于 $I_X$ 的\trAlmostComplex{}结构，以及\trSupport{}含于 $\mathrm{int}(X^{\phi,c})$ 之内的\trHamilton{}函数（\trInhomogeneous{}项）。\trMaximumPrinciple{}表明，\trInhomogeneous{}\trPseudoholomorphic\trPolygon{}保持在 $X^{\phi,c}$ 之内，这意味着 $\scr{F}(X^{\phi,c})$ 作为 $\scr{F}(X^{\phi,d})$ 的一个全 \trAinfty-子\trCategory{}包含于其中。另一方面，$\scr{F}(X^{\phi,d})$ 的每个\trExact{}\trLagrange\trSubmanifold{}都可以经负\trLiouville{}流推入 $X^{\phi,c}$ 之内——在 \trStein{}情形，这与 $\phi$ 的负\trGradient{}流相同。这表明包含 $\scr{F}(X^{\phi,c})\to\scr{F}(X^{\phi,d})$ 是一个\trQuasiequivalence。

其次，\emph{若 $\phi=\psi$ 在一个\trCompact{}子集之外成立，则相伴的\trFukayaCategory{}相同。}可以在两个函数之间\trInterpolation{}并应用一次标准的 \trMoser{}型论证，来证明适当大的\trSublevelSet{}是辛\trIsomorphic{}的，其间是\trExact{}的、并且在边界附近等于恒同的\trSymplectic{}\trIsomorphism。

第三，\emph{假设我们按\cref{thm:plurisubharmonic-monotone-recomposition}从 $\phi$ 过渡到某个 $h(\phi)$。那么\trFukayaCategory{}保持不变。}我们可以选一个辅助函数 $\tilde{h}$，它具有相同的性质，在 $r\gg 0$ 时与 $h$ 一致，并且对 $r\le c$ 有 $\tilde{h}(r)=ar+b$，其中常数 $a>0$、$b\in\RR$。从 $\phi$ 过渡到 $\tilde{h}(\phi)$ 只以一次常数\trRescaling{}改变了\trSublevelSet{}$X^{\phi,c}$ 的辛几何，这可以经调整\trFloerDatum{}与\trPerturbationDatum{}轻易纳入考量。由于两个函数在该\trSublevelSet{}之外都没有\trCriticalPoint，相伴的\trFukayaCategory{}是一致的。另一方面，从 $\tilde{h}(\phi)$ 过渡到 $h(\phi)$ 是一个\trCompactlySupported{}的改变，于是上面给出的论证适用。

经过这些准备之后，我们来讨论主论证。设给定 $\phi$、$\psi$，它们是\trExhausting{}的\trFiniteType{}\trPlurisubharmonic{}函数。由引理 19.1 可找到某个函数 $h(\psi)$，其增长速度快于 $\psi$。取某个 $d$ 大于 $\psi$ 的一切\trCriticalValue。引理 19.2 随之给出一个 $\tilde{\phi}$，它在无穷远处等于 $h(\psi)+(\text{常数})$，并且在 $X^{\psi,d}$ 上与 $\psi$ 一致。现在再运行一遍相同的论证：取某个 $c$ 大于 $\tilde{\phi}$ 的一切\trCriticalValue，且使得 $X^{\psi,d}\subset\mathrm{int}(X^{\tilde{\phi},c})$。再次由引理 19.1 可找到某个 $k(\psi)$，它在 $X^{\psi,d}$ 上等于 $\psi$，且增长速度快于 $\tilde{\phi}$。再运用引理 19.2 找出另一个函数 $\tilde{\psi}$，它在无穷远处等于 $k(\psi)+(\text{常数})$，而在 $X^{\tilde{\phi},c}$ 上等于 $\tilde{\phi}$。

由上面的讨论可知，与 $\tilde{\phi}$ 相关联的\trFukayaCategory{}与 $\phi$ 的那个相同。$\tilde{\psi}$ 与 $\psi$ 亦然；但在后一种情形，我们还有一个额外的观察。考虑 $(X,\phi_{\tilde{\psi}},\theta_{\tilde{\psi}})$ 与 $(X,\phi_{k(\psi)},\theta_{k(\psi)})$ 之间的\trSymplectic{}\trIsomorphism，它由在两个函数之间\trInterpolation{}并应用 \trMoser{}技巧而得。这个\trIsomorphism{}在 $X^{\psi,d}$ 上是恒同的，那里 $\tilde{\psi}$、$\tilde{\phi}$、$\psi$ 和 $k(\psi)$ 在相差加常数的意义下全部一致。$k(\psi)$ 的负 \trLiouville{}流会把 $(X,\phi_{k(\psi)},\theta_{k(\psi)})$ 的每个\trCompact{}子集推入该\trSublevelSet。于是，应用 \trMoser{}辛\trIsomorphism{}之后，$(X,\omega_\psi,\theta_\psi)$ 的每个\trExact{}\trLagrange{}\trSubmanifold{}都\trIsotopy{}于位于 $X^{\psi,d}$ 内部的某一个。

取 $\tilde{d}\gg c$。包含 $(X^{\tilde{\phi},c},\phi_{\tilde{\phi}},\theta_{\tilde{\phi}})\hookrightarrow(X^{\tilde{\psi},\tilde{d}},\phi_{\tilde{\psi}},\theta_{\tilde{\psi}})$ 按命题 10.5 的方式在\trFukayaCategory{}上诱导一个\trCohomology{}\trFullyFaithful{}的\trFunctor。另一方面，该包含的像包含 $X^{\psi,d}$。因此，由前面的观察可知该\trFunctor{}是一个\trQuasiequivalence。
\end{proof}
\begin{remark}\label{thm:polyball-alternative}
还有把一个 \trStein{}\trManifold{}裁成\trCompact{}子集的其他可能办法。例如，若 $X\subset\CC^N$ 是一个\trAffineVariety，可以考虑它与一个\textbf{\trPolyball}（polyball）
\begin{equation*}
P=\{|x_1|^2+\cdots+|x_{N_1}|^2\le c_1,\ \ldots,\ |x_{N_1+\cdots+N_{r-1}+1}|^2+\cdots+|x_N|^2\le c_r\},
\end{equation*}
的交 $M=X\cap P$，其中 $N_1+\cdots+N_r=N$，而 $c_1,\ldots,c_r$ 充分大。选取一个\trPlurisubharmonic{}函数，使它的负\trLiouville{}流关于任何变量 $x_k$ 都指向内部（例如标准函数 $\tfrac14|x|^2$，或 Fubini--Study \trKaehlerPotential{}$\log(1+|x|^2)$）。于是，相应 1-\trForm{}与 2-\trForm{}的限制使 $P$ 成为一个\trExactSymplectic\trCorneredBoundary\trManifold。作为\cref{thm:contact-type-embedding-quasi-iso}的推论，\trFukayaCategory{}$\scr{F}(M)$ 将与我们先前那个构造所得的相同。

作为又一种\trEquivalence{}的进路，人们可以忘掉\trCompactness，而使用整个 $X$，采用在无穷远处等于 $I_X$ 的\trAlmostComplex{}结构。这并不十分契合本书的框架，但它或许是最自然的想法，而且也适用于无限型的情形。
\end{remark}
